\chapter{Границы модели и что дальше}\label{ch:limits}

\section{Как правильно сравнивать с опытом}

Модель --- функция: \emph{условия опыта $\to$ структура}. На входе --- сплав (линии
растворимости, модули, предел текучести), текстура и зерно, водород, $T_{\max}$, скорость
охлаждения, нагрузка. На выходе --- пластинки и их меры. Отсюда три правила.
\begin{enumerate}
\item \textbf{Калибровка --- на одном наборе данных.} Подогнанные параметры ($B$, $\Delta_0$, фора
  зёрен) подбираются по одному, самому полному опыту, близкому к целевому материалу.
\item \textbf{Проверка --- по закономерностям.} Остальные опыты моделируются со своими условиями,
  и сравниваются знак и величина эффектов: водорода, $T_{\max}$, нагрузки, скорости охлаждения,
  размера зерна.
\item \textbf{Допуск --- разброс опыта.} Разные лаборатории для «того же» Zircaloy-4 дают пороги от
  75 до 200~МПа; требовать от модели большей точности, чем согласие опытов между собой,
  бессмысленно.
\end{enumerate}

\fig{F24_Params}{Откуда параметры модели. Чем правее колонка, тем слабее обоснование и тем
важнее независимая проверка. Ореолы убрали из правой колонки потолок $\sigma_{\text{cap}}$.}

\section{Реестр допущений}

\begin{table}[htbp]\centering\small
\caption{Главные допущения: что они исключают и как их проверить.}\label{tab:assump}
\begin{tabularx}{\linewidth}{>{\raggedright}p{3.6cm}>{\raggedright}X>{\raggedright\arraybackslash}p{4.2cm}}\toprule
допущение & что исключает & проверка\\\midrule
двумерное сечение $r$--$\theta$, пластинки бесконечны по оси & наклон пластинок из сечения, двухосность, 3D-связность & 3D-движок; продольные шлифы\\
изотропная матрица (Мизес) & анизотропию зёрен, двойникование & кристаллическая пластичность\\
одна ориентация габитуса на зерно & варианты $\{10\bar17\}$, $\{10\bar11\}$ & EBSD гидридов\\
фора зёрен --- подогнанное число & её зависимость от состояния металла & расчёт остаточных напряжений; нейтронография\\
ореолы складываются, $\sigma_y$ постоянен & взаимодействие пластических зон, рост $\sigma_y$ при остывании & стопки на Фурье; таблица от $T$\\
сетка 0.4~мкм при толщине 0.6~мкм & детали поля у кончика & сетка 0.2~мкм\\
поле 240~мкм, периодическое & градиенты по стенке, свободные поверхности & поле во всю стенку\\
зарождение по CNT с $B$, $\Delta_0$ & зарождение на дислокациях, память прежних гидридов & опыты с циклами\\
\bottomrule
\end{tabularx}
\end{table}

\section{Что модель не может описать по построению}

\paragraph{Двухосность в 2D.} У Cinbiz осевое напряжение снижает порог вдвое. Но и радиальная, и
окружная пластинка содержат ось трубы, и их несоответствие вдоль оси одинаково ($\eps_t$). Значит,
$\sigma_z\eps^*_{zz}$ у обеих одинаков, и в любой модели с \emph{изотропной} матрицей (упругой или
пластичной по Мизесу) осевое напряжение выбор между ними не меняет. Мизес к тому же не чувствует
давление. Поэтому двухосность --- не честный тест для этой модели; это открытый вопрос, для
которого нужна анизотропия на масштабе зерна.

\begin{exercise}
Докажите утверждение про двухосность: запишите $\sigma:\eps^*$ для радиальной и окружной
пластинки при $\sigma=\mathrm{diag}(\sigma_\theta,0,\sigma_z)$ и найдите, от чего зависит их разность.
Почему в пластичной матрице Мизеса это остаётся верным? (Подсказка: девиатор, симметрия
поворота на $90^\circ$ вокруг оси трубы.)
\end{exercise}

\section{Как найти естественные пределы модели}

\begin{enumerate}
\item \textbf{Внутренняя неопределённость.} Прогоны с разными затравками и размером поля, сеткой
  0.4 и 0.2~мкм, ореолами и потолком, $\sigma_y\pm20\%$ дают полосу для каждого выхода. Если полоса
  шире разброса опыта, этот выход модель не предсказывает --- и подгонять его нечего.
\item \textbf{Таблица опытов.} Для каждой работы --- материал, текстура, зерно, водород, $T_{\max}$,
  скорость, вид нагрузки, мера RHF, разрешение, число образцов, разброс. Тогда видно, какие
  расхождения модели меньше разброса между работами.
\item \textbf{Различимость параметров.} Какие данные определяют какой параметр: порог ---
  фору; крутизну и время --- $B$, $\Delta_0$, $k_{\text{tip}}$; долю межзёренных --- фору границы.
  Если два параметра двигают выход одинаково, по этим данным их не разделить.
\end{enumerate}

\section{Открытые вопросы}

\begin{itemize}
\item \textbf{Откуда фора зёрен.} Вероятно, остаточные межзёренные напряжения после прокатки.
  Проверка --- расчёт текстурованного поликристалла с кристаллической пластичностью (EPSC или
  Фурье) или нейтронографические данные. Тот же расчёт может объяснить двухосность: нагрузка по-разному
  делится между зёрнами при одноосном и двухосном растяжении.
\item \textbf{Модель с ореолами.} Перекалибровка, проверка сетки 0.2~мкм, таблица для Э635,
  затем Son (ширина перехода, доля межзёренных) и связность.
\item \textbf{Скорость охлаждения.} Гипотеза: под нагрузкой радиальные стопки побеждают ростом, а
  не первым зарождением; тогда медленное охлаждение даёт им время собрать водород.
\item \textbf{3D.} Шлиф --- сечение; трёхмерная модель нужна для продольных сечений и связности.
\end{itemize}

\begin{exercise}
Придумайте опыт, который различит две гипотезы о природе форы зёрен: (а) остаточные
напряжения после прокатки, (б) плотность дислокаций. Что изменится после отжига, снимающего
напряжения, но не рекристаллизующего металл?
\end{exercise}

\begin{idea}[Итог]
Модель держится на трёх простых идеях: (1) гидрид выпадает, когда химическая движущая сила вместе
с механикой перевесит барьер зарождения, и делает это лавиной; (2) кто первым начал, тот забирает
водород, поэтому порог = фора / 0.08; (3) пластинки взаимодействуют через поле напряжений,
сглаженное пластичностью, и складываются в цепочки и стопки. Всё остальное --- детали, и в
каждой детали есть допущение, которое можно проверить.
\end{idea}
